\documentclass[10pt]{article}
\usepackage[margin=0.72in]{geometry}
\usepackage{booktabs}
\usepackage{tabularx}
\usepackage{array}
\usepackage{enumitem}
\usepackage{graphicx}
\usepackage{float}
\usepackage[hidelinks]{hyperref}
\setlength{\parindent}{1.2em}
\setlength{\parskip}{0.35em}
\setlist[itemize]{leftmargin=1.2em,itemsep=0.1em,topsep=0.15em}
\renewcommand{\arraystretch}{1.05}
\usepackage[T1]{fontenc}
\usepackage{lmodern}
\usepackage{amsmath}
\usepackage{tikz}
\usetikzlibrary{arrows.meta,positioning,fit,calc}

\title{\vspace{-1.0cm}
ECE 475 Term Project: Branch-Predictor Design Exploration on an
Out-of-Order RISC-V Processor\\
}
\author{%
  Wonju Lee \and Shivam Kak\\
}
\date{}

\begin{document}
\maketitle

\begin{abstract}
We add a configurable front-end branch predictor to the lab's
single-issue out-of-order RISC-V processor (\texttt{riscvooo}, the
Lab~4 deliverable with reorder buffer and scoreboard) and evaluate
three predictor algorithms of increasing complexity: a 1-bit BHT,
a 2-bit saturating-counter BHT, and a two-level adaptive predictor
that XORs the program counter with an 8-bit global history register
to index a single shared pattern history table (the GShare variant
from McFarling 1993). On top of these, we implement two stretch
goals from the proposal --- an 8-deep return-address stack (RAS) and
F-stage JAL prediction --- gated behind a single
\texttt{BP\_ENABLED} define so the unmodified path is byte-equivalent
to lab4. Across the four lab~4 ubmarks, the 2-bit BHT lifts mean IPC
from a baseline of 0.741 to 0.786 ($+6.1\%$), and adding JAL
prediction takes that to 0.799 ($+7.8\%$), which is the winning
configuration. The two-level predictor consistently underperforms
BHT-2 on this workload because the kernels are dominated by
single-direction loop branches that pollute the GHR. The RAS is
functionally correct but is exercised only once or twice per kernel.
Baseline \texttt{*-ooo.out} files are byte-identical to the lab4
reference outputs.
\end{abstract}

\section{Design}

\subsection{Overview and proposal scope}

The original proposal targeted ``the dual-issue in-order superscalar
RISC-V processor we have designed for Lab~3,'' but the actual lab tree
contains a single-issue in-order seven-stage pipeline
(\texttt{riscvlong}) and a single-issue out-of-order pipeline with
reorder buffer (\texttt{riscvooo}, the Lab~4 deliverable). This
project targets \texttt{riscvooo} so the predictor is integrated with
the actual term-of-record OoO core, which means the misprediction
penalty is two cycles (squash F + D when a conditional branch is
wrong-path) rather than the three cycles the proposal cited for an
X0-resolution dual-issue design. All three predictor algorithms apply
unchanged. The proposal also called for an FPGA synthesis flow on a
Nexys-4 DDR / Artix-7 part, which is deferred because the lab bench
host became unreachable across the project window (sshd accepted TCP
on port 22 but stopped sending its banner, and without bench access
no one was able to power-cycle it). The writeup below covers IPC
only.

\subsection{Predictor sub-package}

All seven predictor source files live in a new \texttt{bp/}
sub-package that follows lab4's mcppbs convention. The pre-decoder
(\texttt{bp\_predecode.v}) is a pure-combinational decoder over the
F-stage instruction word that classifies the opcode as a B-type
branch, JAL, or JALR; recognises calls and returns
($\texttt{is\_call}=(\texttt{is\_jal}\lor\texttt{is\_jalr})\land
\texttt{rd}{=}x_1$ and
$\texttt{is\_return}=\texttt{is\_jalr}\land\texttt{rs1}{=}x_1\land
\texttt{rd}{\neq}x_1$); and computes both the branch and JAL targets
from F-stage bits alone:
$\texttt{br\_target}=\texttt{pc}+\text{sext}(\texttt{imm\_sb})$ for
B-type and $\texttt{jal\_target}=\texttt{pc}+\text{sext}(\texttt{imm\_uj})$
for JAL. JALR's target requires \texttt{rs1} and is left to the
existing 1-cycle D-stage redirect.

The 1-bit BHT (\texttt{bp\_bht1.v}) is a 256-entry register array
with one bit per entry, indexed by $\texttt{pc}[9{:}2]$ on read and on
write. The 2-bit BHT (\texttt{bp\_bht2.v}) replaces each entry with
the standard \texttt{00}/\texttt{01}/\texttt{10}/\texttt{11}
saturating-counter scheme; the prediction is the upper bit, and on
reset all entries initialise to \texttt{01} (weakly not-taken).

The two-level predictor (\texttt{bp\_two\_level.v}) is a two-level
adaptive scheme in the Yeh \& Patt~1991 family. Level-1 is an 8-bit
global branch history register (GHR) shared across all branches.
Level-2 is a single 256-entry pattern history table (PHT) of 2-bit
saturating counters. The PHT is indexed by
$\texttt{pc}[9{:}2] \oplus \texttt{ghr}[7{:}0]$, the specific
PC-XOR-history variant that McFarling~1993 named ``GShare,'' and the
GHR shifts in the resolved direction of every committed conditional
branch. The GHR is only updated at \emph{resolve} time (X), not at
predict time (F). This avoids needing a checkpoint stack to roll
back the GHR on a mispredict. The cost is that branches in the
F$\rightarrow$D$\rightarrow$X shadow of an as-yet-unresolved branch
see a slightly stale GHR, which is acceptable for the IPC range these
benchmarks reach.

The return address stack (\texttt{bp\_ras.v}) is an 8-entry
$\times$ 32-bit LIFO. On a call, F pushes $\texttt{pc}+4$; on a
return, F pops the top and uses it as the predicted JALR target. The
selectable wrapper (\texttt{bp\_top.v}) instantiates exactly one of
the four direction predictors based on the \texttt{BP\_*} define set
at build time (\texttt{BP\_STATIC\_NT}, \texttt{BP\_BHT1},
\texttt{BP\_BHT2}, or \texttt{BP\_TWO\_LEVEL}). The pre-decoder,
predictor instantiation, and pipeline registers are all gated by
\texttt{`ifdef BP\_ENABLED}, so the unmodified path through
\texttt{riscvooo} is bit-equivalent to lab4 when \texttt{BP\_DEFINES}
is empty.

\subsection{Integration into \texttt{riscvooo}}

The integration patches affect four files:
\texttt{riscvooo-Core.v} (signal wiring),
\texttt{riscvooo-CoreCtrl.v} (predictor instantiation, F-stage
redirect, X-stage resolve, diagnostic counters),
\texttt{riscvooo-CoreDpath.v} (PC-mux extension and
\texttt{pred\_target} pipeline register), and
\texttt{riscvooo-sim.v} (printing the new diagnostic counters in the
output \texttt{*-ooo.out} file). Every change is gated by
\texttt{`ifdef BP\_ENABLED}, so the OoO machinery downstream of
F$\rightarrow$D --- the scoreboard and the reorder buffer we
implemented in Lab 4 --- is untouched. Wrong-path instructions are
squashed in F+D \emph{before} they enter issue, which means the ROB
never sees an instruction that won't commit; no checkpoint stack,
no rollback machinery, and no ROB-flush logic are required. The
predictor is purely a front-end addition.

The pre-decoder runs combinationally on the
\texttt{imemresp\_queue} mux output --- the 32-bit instruction word
that D will receive next cycle --- tagged with \texttt{pc\_Fhl}, and
the predictor is queried with the same PC. The F stage fires a
redirect through a new PC-mux input when the instruction is a
predicted-taken B-type, when it is a JAL (under \texttt{BP\_PRED\_JAL}),
or when it is a JALR-return whose RAS top is valid (under
\texttt{BP\_RAS}). Target priority is RAS, then JAL, then B-type. The
PC-mux ordering becomes (highest priority first)
$\texttt{brj\_taken\_Xhl} \!\to\!
 \texttt{brj\_taken\_Dhl} \!\to\!
 \texttt{pred\_redirect\_Fhl} \!\to\!
 \texttt{pc}+4$, with a new mux input
$\texttt{pm\_pred} = 3$ feeding \texttt{pred\_target\_Fhl} into the PC
register on the next clock edge; the \texttt{pc\_mux\_sel\_Phl}
control signal grows from two bits to three bits when the predictor
is enabled. A small register chain pipelines
$\texttt{predict\_taken}_{\!F}$ along with the \texttt{is\_jal} and
\texttt{ras\_predict} flags from F$\rightarrow$D$\rightarrow$X so
the X-stage compare can determine mispredicts; the bit is valid only
when the X-stage instruction is itself a B-type
($\texttt{br\_sel\_Xhl} \neq \texttt{br\_none}$), and the X-stage
compare gates on \texttt{bp\_resolve\_Xhl} so stray ``predictions''
on non-branches never pollute the predictor's state.

The most important change in the pipeline is
$\texttt{brj\_taken\_Xhl} = \texttt{mispredict\_Xhl}$. With the
predictor on, the X-stage redirect only fires when the prediction was
wrong: a correct taken-prediction does not squash F+D because F is
already fetching from the correct target, and a correct
not-taken-prediction also does not squash. Correct predictions cost
zero cycles, and that is where the IPC improvement comes from. On a
mispredict the redirect target switches between \texttt{branch\_targ}
(actual taken) and $\texttt{pc\_Xhl}+4$ (actual not-taken) via a new
\texttt{redirect\_to\_target\_Xhl} signal back to the dpath. Without
prediction, JAL and JALR raise \texttt{brj\_taken\_Dhl} to redirect
F in D (the existing 1-cycle penalty); under JAL prediction the
redirect is suppressed when \texttt{is\_jal\_Dhl} matches the
pipelined \texttt{pred\_jal\_Dhl} flag, and on a JALR-return whose
RAS prediction matches the actual register-computed target a 32-bit
equality check in dpath (\texttt{pred\_target\_match\_Dhl}) gates the
D-stage redirect off. Wrong RAS predictions still cost the existing
1-cycle penalty.

\section{Methodology}

We deliberately constrained the design so that the build flow,
toolchain, and output format match lab4's exactly. With
\texttt{BP\_DEFINES} unset, the project produces \texttt{*-ooo.out}
files that are byte-identical to lab4's reference outputs for all
47 asm tests and all four ubmarks; only the trailing
\texttt{\$finish} line-number annotation differs, because
\texttt{riscvooo-sim.v} grew seven lines for the
\texttt{`ifdef BP\_ENABLED} diagnostic prints. Predictor variants
are selected by passing one or more \texttt{BP\_*} defines at
\texttt{make} time:
\begin{verbatim}
make BP_DEFINES="-DBP_ENABLED -DBP_BHT2 -DBP_PRED_JAL -DBP_RAS" \
     riscvooo-sim check-asm-riscvooo run-bmark-riscvooo
\end{verbatim}

The cycle and instruction counts are read from
\texttt{proc.ctrl.num\_cycles} and \texttt{proc.ctrl.num\_inst} ---
the same registers \texttt{riscvooo-sim.v} reports in lab4 --- gated
by $\texttt{stats\_en} \lor \texttt{csr\_stats}$. The ubmark
\texttt{make} target passes \texttt{+verbose=1} (no \texttt{+stats}),
so the counters fire only inside the \texttt{test\_stats\_on} /
\texttt{test\_stats\_off} kernel region of each benchmark, which is
the convention lab4 uses; the asm-test rule passes \texttt{+stats=1}
and counts the whole program. Five new \texttt{BP\_ENABLED}-gated
counters --- \texttt{num\_branches}, \texttt{num\_taken},
\texttt{num\_pred\_resolve}, \texttt{num\_pred\_correct},
\texttt{num\_mispredicts} --- share the same gating clause.

We sweep nine BP variants: baseline, \texttt{static\_nt}, BHT-1,
BHT-2, two-level, plus three pairings of BHT-2 and the two-level
predictor with JAL and JAL+RAS. Each variant is a clean rebuild
($\sim$10\,s) followed by all 47 asm tests and 4 ubmarks
($\sim$5\,s), and the full sweep takes about 90 seconds on
Princeton's \texttt{adroit-vis} cluster. All 47 asm-test targets in
\texttt{l4/lab4/build/Makefile} pass on all nine variants
($47 \times 9 = 423$ runs, all \texttt{[ PASSED ]}). This includes
the standard branch-and-jump suite (\texttt{riscv-beq},
\texttt{riscv-bne}, \texttt{riscv-blt}, \texttt{riscv-bge},
\texttt{riscv-bltu}, \texttt{riscv-bgeu}, \texttt{riscv-j},
\texttt{riscv-jal}, \texttt{riscv-jalr}, \texttt{riscv-jr}), the
integer/load-store/mul-div suite, and the four OoO-specific tests we
wrote in Lab 4 (\texttt{riscv-ooo-commit}, \texttt{riscv-rob-bypass},
\texttt{riscv-waw}, \texttt{riscv-long-better-ipc}). All four ubmarks
(\texttt{vvadd}, \texttt{cmplx-mult}, \texttt{bin-search},
\texttt{masked-filter}) report \texttt{*** PASSED ***} on every
variant. The predictor integration preserves architectural
correctness in every configuration we tested.

\section{Evaluation}

Tables~\ref{tab:ipc_main} and~\ref{tab:ipc_full} report the
kernel-only IPC for the four ubmarks across all nine variants, and
Table~\ref{tab:misp} reports the per-variant mispredict counts on the
$N$ kernel-region conditional branches.

\begin{table}[h]
\centering
\caption{IPC by variant (kernel-only, lab4 convention). Five-variant
view; \texttt{static\_nt} is the integration-overhead control.}
\label{tab:ipc_main}
\begin{tabular}{lrrrrr}
\toprule
Benchmark & baseline & static-NT & BHT-1 & BHT-2 & two-level (PC$\oplus$GHR) \\
\midrule
\texttt{vvadd}        & 0.8865 & 0.8865 & 0.9115 & 0.9115 & 0.8830 \\
\texttt{cmplx-mult}   & 0.7105 & 0.7105 & 0.7236 & 0.7236 & 0.7194 \\
\texttt{bin-search}   & 0.7048 & 0.7048 & 0.7865 & 0.7952 & 0.7664 \\
\texttt{masked-filter}& 0.6622 & 0.6622 & 0.7064 & 0.7153 & 0.7115 \\
\midrule
\textbf{mean}         & \textbf{0.741} & \textbf{0.741} & \textbf{0.782} & \textbf{0.786} & \textbf{0.770} \\
\bottomrule
\end{tabular}
\end{table}

\begin{table}[h]
\centering
\caption{IPC with the proposal stretch goals (JAL prediction at F,
8-deep RAS for JALR returns) layered on top of BHT-2 and the
two-level predictor. ``full'' = direction predictor + \texttt{BP\_PRED\_JAL}
+ \texttt{BP\_RAS}.}
\label{tab:ipc_full}
\begin{tabular}{lrrrr}
\toprule
Benchmark & BHT-2 + JAL & two-level + JAL & BHT-2 full & two-level full \\
\midrule
\texttt{vvadd}        & 0.9115 & 0.8830 & 0.9115 & 0.8830 \\
\texttt{cmplx-mult}   & 0.7239 & 0.7197 & 0.7239 & 0.7197 \\
\texttt{bin-search}   & 0.8215 & 0.7908 & 0.8215 & 0.7908 \\
\texttt{masked-filter}& 0.7394 & 0.7354 & 0.7394 & 0.7354 \\
\midrule
\textbf{mean}         & \textbf{0.799} & \textbf{0.782} & \textbf{0.799} & \textbf{0.782} \\
\bottomrule
\end{tabular}
\end{table}

\begin{table}[h]
\centering
\caption{Conditional branches in the kernel region and mispredicts
per direction predictor. Percentages are mispredicts/branches.}
\label{tab:misp}
\begin{tabular}{lrrrr}
\toprule
Benchmark & branches & BHT-1 & BHT-2 & two-level (PC$\oplus$GHR) \\
\midrule
\texttt{vvadd}         &   10 &   2 (20.0\%) &   2 (20.0\%) &  10 (100.0\%) \\
\texttt{cmplx-mult}    &   27 &   3 (11.1\%) &   3 (11.1\%) &  10 ( 37.0\%) \\
\texttt{bin-search}    &  246 &  59 (24.0\%) &  52 (21.1\%) &  76 ( 30.9\%) \\
\texttt{masked-filter} &  668 &  96 (14.4\%) &  53 ( 7.9\%) &  71 ( 10.6\%) \\
\bottomrule
\end{tabular}
\end{table}

As a sanity check against lab4, the baseline \texttt{*-ooo.out} files
diff cleanly against \texttt{l4/lab4/build/*-ooo.out}: $511/453/0.886$
on \texttt{vvadd}, $2425/1723/0.711$ on \texttt{cmplx-mult},
$1443/1017/0.705$ on \texttt{bin-search}, and $7446/4931/0.662$ on
\texttt{masked-filter} (\texttt{cycles}/\texttt{inst}/\texttt{ipc}).
The build flow and RTL behave bit-for-bit like lab4's reference for
the unmodified core.

\subsection{Discussion}

The first thing to note is that \texttt{baseline} and
\texttt{static\_nt} produce identical IPC to four decimal places on
every benchmark. The lab core's existing behaviour is itself
``always not-taken'' --- fetch PC+4, correct on X resolution --- so
adding the predictor in \texttt{static\_nt} mode should not change
observable behaviour, and it does not. This validates that the new
pipeline registers, the widened PC mux, and the \texttt{`ifdef BP\_ENABLED}
hooks introduce no correctness regression or measurable overhead
when the predictor signal is constant. Among the actually-predicting
configurations, BHT-1 and BHT-2 tie on three of four benchmarks; the
exception is \texttt{masked-filter}, where BHT-2 cuts mispredicts
almost in half (96 $\to$ 53, $-45\%$). That benchmark contains an
inner branch whose direction alternates at a frequency that exposes
BHT-1's weakness: a single anomalous outcome flips the prediction, so
a nearly-deterministic mostly-taken branch with one rare not-taken
iteration costs four mispredicts (one to flip out, one for the rare
case, one to flip back, one for the next normal case). The 2-bit
hysteresis absorbs the rare case without flipping the prediction.

The two-level predictor underperforms BHT-2 on every benchmark, and
the most striking case is \texttt{vvadd}, where every single one of
the ten kernel-region conditional branches is mispredicted. The
underlying reason is GHR pollution: the kernel's branches are nearly
all single-direction loop branches with no useful global-history
correlation, but the GHR cycles through patterns of mostly-1s
anyway, so the same PC maps to different PHT counters depending on
how many other taken branches resolved recently. Each mapped counter
then has only a handful of training cycles to bias toward ``taken''
before it gets remapped, so the warm-up cost dominates the kernel.
BHT-2's per-PC counters do not suffer this --- they accumulate
hysteresis across all visits to the same branch. The general lesson
is that the global-history correlation that the two-level scheme is
built around only pays off when there is correlation to exploit; on
\texttt{ubmark-vvadd}/\texttt{cmplx-mult}/\texttt{bin-search}/
\texttt{masked-filter}, there is none.

The cycle accounting on \texttt{vvadd} makes the small numbers in
the table concrete. With \texttt{-O3 -funroll-loops} (matching lab4's
\texttt{ubmark.mk}), the 100-iteration loop collapses into roughly
12 unrolled chunks; only ten conditional branches survive in the
measured region, of which nine are taken. The baseline kernel runs
in 511 cycles and retires 453 instructions (IPC $0.886$). The
static cost of mispredicting nine taken branches is
$9 \times 2 = 18$ cycles. BHT-2 cuts that to two mispredicts
$\times$ two cycles $= 4$ cycles, so the kernel completes in 497
cycles (IPC $0.911$): a saving of 14 cycles, $\approx 78\%$ of the
18-cycle ceiling. The remaining gap is from cold predictor entries
on the very first loop iteration plus a couple of mispredicts in
the chunk-out tail where the unroll factor doesn't divide the
iteration count evenly.

JAL prediction and the RAS together lift mean IPC from BHT-2's
$0.786$ to $0.799$, but the contribution is uneven across
benchmarks and across the two stretch goals. JAL prediction is the
larger contributor: on \texttt{bin-search} it lifts IPC from $0.795$
to $0.822$ ($+3.4\%$), and on \texttt{masked-filter} from $0.715$ to
$0.739$ ($+3.4\%$). \texttt{bin-search} contains the inner-loop call
sites where this stretch goal pays off, while on \texttt{vvadd} and
\texttt{cmplx-mult} the kernel JAL count is zero or one and the IPC
delta is below counter resolution. The RAS, by contrast, is
functionally correct but adds zero IPC on these benchmarks ---
\texttt{bp\_bht2\_full} produces the same numbers as
\texttt{bp\_bht2\_jal} on every benchmark. The reason is that the
\texttt{test\_stats\_on/off} kernel regions of these ubmarks contain
at most a couple of JALR returns (most return-from-main happens
\emph{outside} \texttt{test\_stats\_off}), so on a workload with
deeper recursion or more function calls inside the measured region
the RAS would matter more. We verified RAS correctness in waveform:
the \texttt{push\_addr} on a JAL-with-rd=$x_1$ matches the next
sequential PC, the \texttt{pop} on a JALR-with-rs1=$x_1$/rd=$\neq x_1$
recovers that PC, and \texttt{pred\_target\_match\_Dhl} suppresses
the D-stage redirect.

Putting these observations together, BHT-2 + JAL is the highest-IPC
configuration across all nine design points at $0.799$ mean, and
BHT-1 is the smallest predictor that gets within $5\%$ of it
($0.782$ mean). The two-level predictor costs more state than BHT-2
(an extra 8-bit GHR plus an XOR network on the index path) and
produces strictly worse IPC on these workloads, so for this lab's
benchmark set additional history correlation does not pay off; a
workload with strongly correlated inter-branch behaviour would be
expected to flip that conclusion. On a separate axis, we re-ran the
same predictor variants on \texttt{riscvlong} (single-issue
in-order, the Lab~2 deliverable) for comparison and observed two
patterns: baseline IPC is consistently lower on \texttt{riscvooo}
than on \texttt{riscvlong} for these workloads, in line with our
Lab~4 report's conclusion that tight dependency chains expose the
OoO front-end's extra pipeline state and in-order commit constraint
without recovering the cost; and the \emph{relative} IPC lift from
the predictor is similar across the two cores, because the
misprediction penalty is two cycles in both pipelines and the cycles
the predictor saves are the same cycles in both.

\section{Related work}

The 1-bit and 2-bit BHT designs follow Smith~\cite{smith1981}, who
introduced the saturating-counter scheme that all subsequent direction
predictors build on. The two-level adaptive structure --- a global
history register feeding a pattern history table --- is the
Yeh \& Patt~\cite{yehpatt1991,yehpatt1992} family, and the specific
PC$\oplus$GHR indexing variant we implemented is from
McFarling~\cite{mcfarling1993}, who also introduced tournament
predictors that combine local and global predictors and showed
GShare's accuracy benefit on SPEC at small storage budgets. The
8-deep return-address stack follows Kaeli \&
Emma~\cite{kaeliemma1991}, who first proposed using a stack to
predict return-from-call targets in pipelined processors, and
Hennessy \& Patterson~\cite{hp6e} surveys the trade-offs at textbook
level. More recent academic work, particularly perceptron
predictors~\cite{jimenez2001} and TAGE~\cite{seznec2006tage},
substantially outperforms the predictors in this study on full-system
workloads, but those designs require considerably more state and
several pipelined stages of lookup, and they would not fit cleanly
into the existing single-cycle F-stage of \texttt{riscvooo} without
re-pipelining. The proposal originally called for an FPGA area/Fmax
tradeoff in the spirit of McFarling and several follow-on FPGA
branch-predictor studies; that analysis is the obvious next step for
a future revision of this work once a synthesis bench is available.

\section{Conclusion}

This project added a configurable front-end branch predictor to the
Lab~4 OoO core, evaluated three direction-predictor algorithms plus
two stretch goals (JAL prediction and a return-address stack), and
ran the four lab~4 ubmarks through all nine variants. The
\texttt{BHT-2 + JAL} configuration is the winning point at $0.799$
mean IPC ($+7.8\%$ over baseline). The two-level predictor with
PC$\oplus$GHR-indexed PHT consistently underperforms the simpler
2-bit BHT on this workload because the kernels contain many
single-direction loop branches that pollute the GHR. The 8-deep RAS
is functionally correct but the kernels exercise it only minimally;
its mechanism is in place for workloads that contain deeper
recursion or more function calls. The implementation lives in a new
\texttt{bp/} sub-package and four lightly patched files in
\texttt{riscvooo/}, the build flow uses lab4's existing autoconf
pipeline plus one new \texttt{BP\_DEFINES} knob, and the baseline
\texttt{*-ooo.out} files are byte-identical to the lab4 reference.
The full sweep is reproducible from a clean checkout with
\texttt{source scripts/adroit-env.sh \&\& ./scripts/run\_variants.sh},
and the source is at
\url{https://github.com/Shivamkak19/fpga_riscv_branch_predictor}
on the \texttt{main} branch.

\begin{thebibliography}{9}

\bibitem{smith1981}
J.~E.~Smith.
\emph{A study of branch prediction strategies.}
ISCA, 1981.

\bibitem{yehpatt1991}
T.-Y.~Yeh and Y.~N.~Patt.
\emph{Two-level adaptive training branch prediction.}
MICRO, 1991.

\bibitem{yehpatt1992}
T.-Y.~Yeh and Y.~N.~Patt.
\emph{Alternative implementations of two-level adaptive branch
prediction.}
ISCA, 1992.

\bibitem{mcfarling1993}
S.~McFarling.
\emph{Combining branch predictors.}
WRL Technical Note TN-36, 1993.

\bibitem{kaeliemma1991}
D.~R.~Kaeli and P.~G.~Emma.
\emph{Branch history table prediction of moving target branches due
to subroutine returns.}
ISCA, 1991.

\bibitem{jimenez2001}
D.~A.~Jim\'enez and C.~Lin.
\emph{Dynamic branch prediction with perceptrons.}
HPCA, 2001.

\bibitem{seznec2006tage}
A.~Seznec and P.~Michaud.
\emph{A case for (partially) tagged geometric history length branch
prediction.}
JILP, 2006.

\bibitem{hp6e}
J.~L.~Hennessy and D.~A.~Patterson.
\emph{Computer Architecture: A Quantitative Approach}, 6th ed.
Morgan Kaufmann, 2017. Chapter on branch prediction.

\end{thebibliography}

\end{document}
